\subsubsection{The effective cohomology bound}\label{tw:effectiveK}

Let $p$ be a prime, $L$ a totally imaginary number field of degree $d$ and
absolute discriminant $D$, and $S$ the set of places above $p$.
Suppose $L$ contains a primitive $p$th root of unity if $p$ is odd
and a primitive fourth root if $p=2$. Write
\[
 B_C(e,d,D)=\exp((2ed)^C\log(2D)).
\]
The bound required by the arithmetic window argument is uniform: for one
real $C_K$, all such $p,L$ and all integers $e\ge1$, the group
$A=H^2(G_{L,S},\mathbb Z_p(e+1))$ is finite and
$\log |A|\le B_{C_K}(e,d,D)$. In forming the cochains, fix an algebraic
closure, a prolongation of each $v\in S$ to it, and the cyclotomic
character on the quotient unramified outside $S$. The estimate holds for
each such datum.

\paragraph{The localization kernel.}
For the discrete module $M=(\mathbb Q_p/\mathbb Z_p)(-e)$ let
$s_L$ denote the length over $\mathbb Z_p$ of the kernel of
\[
 H^1(G_{L,S},M)\longrightarrow \prod_{v\in S} H^1(G_{L_v},M).
\]
Lengths can initially be infinite. The source identifies this length with
the invariant defined using an absolute Galois group and a fixed algebraic
closure. At the end of the argument one returns from this invariant to
the same local conditions. The two twists involved here are $-e$ in degree
one and $e+1$ in degree two.

\paragraph{Counting in cyclotomic layers.}
Put $k=e-1$. The auxiliary field construction supplies a totally imaginary
field $L'$ with $d'=[L':\mathbb Q]=m d$, maximal cyclotomic depth $c$
($\mu_{p^c}\subset L'$, $\mu_{p^{c+1}}\not\subset L'$), and
\[
 2e+1<\min(p^c,r_2(L')), \quad
 m\le(d+1)(4e+4), \quad p^c\le 2d',\quad
 D'\le D^m(p^c)^{d'},\quad s_L\le s_{L'}.
\]
Here $D'=|\Delta_{L'}|$, $r_2$ counts pairs of complex embeddings, and
$c\ge1$, with $c\ge2$ at $p=2$. The strict inequalities are used before
applying the odd-degree homology estimate.

Indeed, if $c_0$ is the cyclotomic depth of $L$, take the layer
$L'=L(\mu_{p^{c_0+h}})$ with $h$ the least integer for which
$p^h\ge4e+4$. The cyclotomic layer theorem gives $m=p^h$.
Thus $m\le p(4e+4)$, $c\ge c_0+h$, and $r_2(L')=md/2>2e+1$.
The cyclotomic degree over $\mathbb Q$ gives
$p\le d+1$ and $p^c\le 2d'$. To bound the different, let $z$ generate
the layer as a primitive $N$th root, $N=p^{c_0+h}$, and let $P$ be
its relative minimal polynomial. Writing $X^N-1=PQ$ and differentiating
at $z$ gives $N=zP'(z)Q(z)$. The different contains $P'(z)$ and hence
$N$. Its norm divides $N^{d'}$, giving the displayed discriminant bound.
For the length comparison use the tower descent description at levels
$0$ and $h$. These lengths are those of the cokernels of $T-1$ and
$T^{p^h}-1$ for the same twisted operator $T$. The second image lies in
the first by the geometric-sum identity, giving a surjection of
cokernels in the required direction.

In a cyclic Galois layer $F/L'$ let $\zeta$ have order $q=p^\nu$,
$\nu\ge c$, and fix a generator $\gamma$ of order $p^{\nu-c}$
with $\gamma\zeta=\zeta^{1+p^c}$. Set $a=(1+p^c)^e$. Count the
functions $\varphi$ from nonzero prime ideals of $\mathcal O_F$ to
$\mathbb Z/q$ satisfying
\[
 \sum_P\operatorname{ord}_P(\alpha)\varphi(P)=0\quad(\alpha\in F^\times),
 \qquad \varphi(\gamma^{-1}P)=a\varphi(P).
\]
The first condition makes these characters of the ideal class group.
Write $h_\nu$ for their number, $T_j$ for the order of the torsion
subgroup of $H_j(\mathrm{GL}_{4k+10}(\mathcal O_{L'}),\mathbb Z)$, and
\begin{align*}
 b&=2k+2+(2k+3)(c+2k+3), & f&=v_p(e),\\
 g&=2d'+\lfloor\log_2 h_{L'}\rfloor+2, & X&=(b+f)g+c+f .
\end{align*}

Two counts enter. For the first count, consider a finite subset of
$F^\times$ whose elements have
$\operatorname{ord}_P(\alpha)=\mathfrak d(\gamma P)-a\mathfrak d(P)$
for a divisor $\mathfrak d$ depending on $\alpha$. The chosen elements
are distinguished modulo $q$th powers by the $p^b$th powers of their
twisted norms. In the archive convention the norm is
$\prod_{j=0}^{|\gamma|-1}\gamma^j(\alpha)^{a^j}$, with the field
automorphism acting on $\alpha$. Each element gives a cycle
in degree $2k+3$. One chooses the residue tests simultaneously away
from $p$ and the finitely many denominators of these cycles, at primes
where $\mathcal O_{L'}$ surjects onto the residue field. If two cycles
have the same scaled evaluations at every such prime, the ratio of their
powered twisted norms is a $q$th power in every tested residue field.
The power-detection lemma gives a $q$th power in $F^\times$.
The separation condition then identifies the two elements. The image
bound for the resulting primitive family of evaluations bounds the first
count by $q^{r_2(L')}T_{2k+3}T_{2k+2}$.

\paragraph{The cycle and its determinant.}\label{kone:cycle}
Write $A_q=\mathbb Z/q$, $C=\mathbb Z/q$ as an additive group,
$M=2k+3$, $\kappa=1+p^c$ and $a=\kappa^e$, so $e=k+1$.
In the bar complex of $C$ with coefficients in $A_q$ use
\[
 z=\sum_{j_1,\ldots,j_e\in C}[j_1|1|\cdots|j_e|1].
\]
This has degree $2e$. In its boundary the first face sums to zero
because its multiplicity is $q$. The remaining faces cancel in pairs
after shifting a summation variable by one. Put $d_t(x)=(x,tx)$,
$0\le t\le M$. The product-group homology in degree $2e$ with these
coefficients has a generating set of size at most $2e+1=M$ (use the
first projection onto $C$). More explicitly, constants, cyclic
shift minus identity and summation give an exact sequence
from the trivial module through two copies of the cyclic permutation
module to the trivial module. The two connecting sequences and
induction from the kernel bound the generator number in degree $d+2$
by that in degree $d$ plus the kernel-group bounds in degrees $d+1,d+2$.
For a cyclic kernel both bounds are one; degree zero needs one
generator. Hence degree $2e$ needs at most $2e+1$ generators.
The $M+1$ classes $(d_t)_*z$
therefore satisfy a relation $\sum_t n_t(d_t)_*z=\partial B$ with
some $n_{t_0}$ a unit. Indeed their subgroup is a finite abelian
$p$-group with at most $M$ cyclic factors. A redundant generator
modulo $p$ is redundant in the subgroup, by successive lifting
through the powers of $p$. This gives the unit coefficient.

Index the coordinates of $F^{C^2}$ by $j=(j_1,j_2)$.
On coordinate $j$ the diagonal representation $\rho$ of $C^2$
has character $(x,y)\mapsto\zeta^{j_1x+j_2y}$. For $t$ let
$s_t(j)=j_1+tj_2$. The lattice construction uses a finitely supported
integer function $\eta$ on $C^2$. Its sums on the fibers of $s_t$
vanish for $t\ne t_0$. At $t_0$ its pushforward is
\[
 \sum_{i=0}^{M}[X^i]P(X)\,[\kappa^i],\qquad
 P(X)=\prod_{j=0}^{M-1}(X-\kappa^j).
\]
Here $[x]$ denotes a basis vector in $\mathbb Z[C]$. To see the
interpolation, use the product of $[(0,1)]-[(t,0)]$ over $t\ne t_0$
in $\mathbb Z[C^2]$. Each unwanted pushforward kills a factor. At
$t_0$ the product equals a translation times
$\prod_{t\ne t_0}([t_0-t]-1)$. Every $t_0-t$ generates a subgroup
containing $p^c$ in $C$, since $0<|t_0-t|\le M<p^c$. Thus each factor
divides $[p^c]-1$. The desired pushforward is divisible by
$([p^c]-1)^M$: in polynomial coordinates apply
$\prod_{j<M}(U-\kappa^j)$ to $X$, where $Uf(X)=f(X^\kappa)$.
The operator $U-\kappa^j$ carries $(X^{p^c}-1)^j$ times any polynomial
into $(X^{p^c}-1)^{j+1}$ times a polynomial, by factoring
$X^{\kappa p^c}-1$ and reducing the quotient modulo $X^{p^c}-1$.
Lift the quotient in the group ring by $x\mapsto(x,0)$ and undo the
translation. This constructs $\eta$ with both specified properties.

Modify the standard lattice coordinate $j$ by the fractional ideal
with divisor $\eta(j)\mathfrak d$. For each $t$ construct an
$L'$-linear $g_t$ taking the standard lattice to this modified lattice
and commuting with $\rho d_t$. The determinant-one identifications on
blocks use the Steinitz relation for a direct sum of fractional ideals:
one can move the sum of two divisors to one coordinate, and iterate.
Scaling one distinguished coordinate realizes a prescribed determinant
when its divisor is the difference of the two total divisors.
For $t\ne t_0$ all block determinants can be taken as one. At $t_0$
the construction groups coordinates in Galois orbits and leaves the
determinant
\[
 \beta=\prod_i\gamma^{-i}(\alpha)^{[X^i]P_e(X)},\qquad
 P_e(X)=\prod_{0\le j<M,\ j\ne e}(X-\kappa^j)
\]
at character $1$. Indeed the divisor of $\gamma^{-i}\alpha$ at $P'$
is $\mathfrak d(\gamma^{i+1}P')-a\mathfrak d(\gamma^iP')$.
Multiplying this difference by the coefficients of $P_e$ gives
precisely the divisor prescribed by $P=(X-a)P_e$. This explains
both the inverse action in $\beta$ and the missing factor in $P_e$.

Let $H$ be the stabilizer of the standard lattice inside the localized
group containing all $g_t$ and $\rho(C^2)$. One base-field denominator
excludes the support of the divisors and the bad primes of the splitting
map. Set $\rho'=g_{t_0}^{-1}\rho g_{t_0}$ and $h_t=g_{t_0}^{-1}g_t$.
Both representations take values in $H$, and $h_t\in H$ since both
maps $g_t,g_{t_0}$ take the standard lattice to the same lattice.
For the bar prism $P_h$ one has
$\partial P_h+P_h\partial=(\mathrm{conj}_{h^{-1}})_*-1$.
Consequently the chain
\[
 Y=\rho'_*B-\rho_*B+\sum_t n_t P_{h_t}((\rho'd_t)_*z)
\]
is closed in $H$. In fact $h_t^{-1}\rho'd_th_t=\rho d_t$, so the
boundary of the sum is $\rho_*\partial B-\rho'_*\partial B$.
In the localized group $Y$ is homologous to
$\sum_t n_t P_{g_t}((\rho d_t)_*z)$. This follows by applying the
prism identity to $\rho_*B$ and the prism product identity to
$g_t=g_{t_0}h_t$ on each cycle. Multiply $Y$ by $n_{t_0}^{-1}$.
The standard lattice is finite projective over $\mathcal O_{L'}$.
Adding a projective complement with trivial action realizes this cycle
in some $\mathrm{GL}_N(\mathcal O_{L'})$.

At an allowed residue field $k_0$, restriction of scalars splits into
the conjugate blocks indexed by $\sigma\in\mathrm{Gal}(F/L')$.
For a character $v:k_0^\times\to A_q$, the finite-field functional
on a scalar-character block of weight $x$ evaluates a prism for $g$
as $x^e v(\det g)$. On the scalar line this is the product of $e$
cyclic carry cocycles with $v$: each consecutive pair in $z$ sums to
$x$, and insertion in the prism leaves $v(g)$. Transfer over projective
lines gives the additive functional on vector spaces. The line count
for direct sums uses $|k_0|=1$ in $A_q$, since the reduced $\zeta$ has
order $q$. Additivity kills the trivial complement. Blocks of determinant
one give zero. In the surviving block let $\sigma\zeta=\zeta^{\chi(\sigma)}$.
Summing the conjugate blocks gives
\[
 \sum_\sigma \chi(\sigma)^e v(\overline{\sigma\beta})
   =P_e(a)\,v\big(\overline{\operatorname{nrm}(\alpha)}\big).
\]
Indeed replacing $\alpha$ by $\gamma^{-i}\alpha$ multiplies the
weighted sum by $a^i$ after reindexing, and the sum for $\alpha$ is
the evaluation of its twisted norm. The unit coefficient cancels the
normalization of $Y$. For a cyclic residue unit group take $v$ with
kernel the $q$th powers. Each factor $\kappa^e-\kappa^j$ of $P_e(a)$
has valuation $c+v_p(|e-j|)\le c+M$, by lifting the exponent
(with $c\ge2$ at $p=2$). In $A_q$, $P_e(a)$ therefore divides
$p^{M(c+M)}$. Equality of two evaluations scaled by $p^{2e}$
implies that the ratio of their twisted norms raised to
$p^{2e+M(c+M)}=p^b$ is a $q$th power at this prime.

\paragraph{Detecting powers and bounding the image.}\label{kone:rank}
If $\xi\in F^\times$ passes all these residue tests, choose
$x$ with $x^q=\xi$ in an algebraic closure. For each automorphism
$g$ over $F$, the condition $\sigma x=gx$ describes an open coset.
Degree-one Chebotarev gives a Frobenius in this coset outside the finite
exclusion set and the primes dividing $q$. Its residue field has prime
order, so reduction from $\mathcal O_{L'}$ is surjective. In the residue
field $\xi$ has a $q$th root and all the $q$th roots of unity are
present. Hence the residue of $x$ is fixed by Frobenius. The quotient
$\sigma x/x$ is a $q$th root of unity congruent to one. If it were
different from one, its geometric sum of length $q$ would be zero,
giving $q=0$ in the residue field. Thus $gx=x$ for every $g$ and
$x\in F$. Applying this argument to the ratio proves injectivity
for the simultaneous evaluations of the finite set.

For the image bound put $m_1=2e$, and stabilize the finitely many
cycles to rank $n\ge2(m_1+1)+4$. An integral lift $Z$ of a
mod-$q$ cycle has $\partial Z=qV$. The homology class of $V$ is
$q$-torsion in degree $m_1$, independent of the lift. When this class
vanishes one can subtract $q$ times an integral chain from $Z$ to
make it a cycle. There are at most $T_{m_1}$ fibers of this obstruction
by homological stability. Differences in a fiber therefore reduce the
count to evaluations on integral cycles.

Write $\psi_p$ for the $p$-fold block sum on integral homology and
$i$ for stabilization to the same rank. Starting with $z_0=x$ put
$z_{t+1}=(\psi_p)_*z_t-p^{t+2}i_*z_t$. A primitive evaluation
$\phi$ is stable and additive on block sums. Hence
$\phi(z_{t+1})=p(1-p^{t+1})\phi(z_t)$. The factors $1-p^{t+1}$
are units modulo $q$. The scaled evaluations $p^{m_1}\phi(x)$ thus
factor through the integral span of the $z_{m_1}$. Its rank is at
most $r_2(L')$. To obtain this last inequality the source passes to
real cycles and pairs with one odd-degree functional at each infinite
place. A class in the common kernel, if nonzero, survives stabilization
and is separated by a closed cochain. The surjectivity theorem represents
that cochain on cycles by an invariant form. Its weight decomposition
consists of the place functionals and terms of weights $j=2,\ldots,m_1+1$.
The latter vanish on the displayed span: the factor at $t=j-2$ is
$p^j-p^j=0$. The former vanish by the kernel condition. This contradicts
separation. The real image thus has dimension at most the number of
complex places, giving the integral rank bound.

For any finite set of evaluations use the finitely generated subgroup
spanned by lifts of its elements. Its free part contributes at most
$q^{r_2(L')}$ values. Its torsion contributes at most $T_{m_1+1}$
by stability, also at the enlarged rank for $z_{m_1}$. Multiplication
by the number of obstruction fibers gives the claimed bound.

\paragraph{Transfer over residue-field lines.}\label{kone:lines}
Here is the transfer used in the cycle calculation of
Paragraph~\ref{kone:cycle}. Let $V$ be a vector space over the finite
residue field $k_0$, let $\rho:G\to\mathrm{GL}(V)$, and let $t$ be the
functional on bars of $k_0^\times$ which vanishes on boundaries.
Choose a representative on each line. This gives a function
$\lambda:V\setminus\{0\}\to k_0^\times$ with
$\lambda(sv)=s\lambda(v)$, by writing $v$ as a multiple of its
representative. Give the lines the right action
$L\cdot g=\rho(g)^{-1}L$. For $0\ne v\in L$, put
\[
 c(g,L)=\lambda(v)/\lambda(\rho(g)^{-1}v).
\]
The ratio is independent of $v$. It satisfies
$c(gh,L)=c(g,L)c(h,L\cdot g)$ by cancellation at the intermediate
line. For a bar $[g_1|\cdots|g_n]$ sum, over initial lines $L_0$,
the bars $[c(g_1,L_0)|\cdots|c(g_n,L_{n-1})]$, where
$L_i=L_{i-1}\cdot g_i$. This sum commutes with the bar boundary.
For the first face reindex by $L_1$; for an interior face use the
cocycle identity; the last face just forgets the last arrow.
Composition with $t$ therefore kills boundaries.

If $\lambda$ is multiplied by a scalar $b(L)$ on line $L$, each
edge label is multiplied by $b(L_{i-1})/b(L_i)$. The two labelled
paths are related by the vertical edges $b(L_i)$: the products along
the two sides of each square agree. Triangulate the prism on each
path by inserting the vertical edge successively at the $n+1$
vertices, with alternating signs. Its interior faces cancel in pairs;
the two ends give the difference of the old and new bars, and the
side faces are the corresponding prisms on the boundary. Thus the
two sums agree under $t$ on cycles. In particular a linear change
of coordinates, which transports both lines and representatives,
preserves the resulting functional.

On $V\oplus W$, send a line to its projection in $V$ when that
projection is nonzero, and otherwise to its line in $W$. Use the
normalizing function of the indicated projection. Each line above
a fixed line in $V$ is the graph of a linear map from that line to
$W$, so its fiber has $|k_0|^{\dim W}$ elements. A line in $W$ has
one preimage. For block diagonal actions the edge labels in a fiber
are exactly those of its projected path. Both fiber sizes equal one
as coefficients in $\mathbb Z/q$, since $q\mid |k_0|-1$. The
functional on a direct sum is consequently the sum of the two
functionals on cycles. On the one-dimensional scalar representation
there is a single line and $c(g,L)=g$, so the construction recovers
$t$. These two facts justify passage from the scalar computation to
the sum over the residue-field blocks.



For the second count write $Y_F=F^\times$ additively, let $I_F$ be the
divisor group, and let $\partial:Y_F\to I_F$ be the divisor map.
On elements put $\delta=\gamma^{-1}-a$; on divisors the corresponding
operator is $(\delta_I d)(P)=d(\gamma P)-a d(P)$. Put
\[
 U=\partial^{-1}(\delta_I I_F),\quad
 W=U/(\delta Y_F+Y_{F,\mathrm{tors}}),\quad
 C'=I_F/(\partial Y_F+\delta_I I_F).
\]
Let $E'$ be $\ker\partial$ modulo $\delta(\ker\partial)$ and torsion.
The divisor and unit counts give finite groups with $|W|=|E'||C'|$,
$h_\nu\le p^{v_p(|C'|)}$, and
\[
 v_p(|E'|)+c+f=r_2(L')(\nu+f).
\]
The last equality uses $v_p(a-1)=c+f$ and
$v_p(a^{p^{\nu-c}}-1)=\nu+f$.

Here are the two order computations. On $I_F$ the operator $\delta_I$
is injective, since an element in its kernel satisfies
$(a^{|\gamma|}-1)d=0$ and the divisor group is torsion free.
For $y\in U$ the equation $\partial y=\delta_I d$ now determines $d$.
Sending $y$ to the ideal class of $d$ maps $W$ onto the kernel of
$\delta_I$ on the class group. Indeed a class in that kernel has
$\delta_I d$ principal, and altering $y$ by $\delta z$ alters $d$ by
$\partial z$. If $d$ is principal one can therefore take $y$ to be
a unit. Two such units agree in $W$ precisely modulo the twisted
differences of units and torsion, again by injectivity on divisors.
Thus the kernel is $E'$. The kernel and cokernel of an endomorphism
of a finite group have the same cardinality, and here the cokernel
is $C'$. Equivariant class characters factor through $C'$; on each
cyclic $p$-summand of order $p^i$ there are at most $p^i$ choices.

The logarithmic unit lattice spans the real hyperplane of coordinates
at infinite places with sum zero. On all coordinates $\gamma^{-1}$
permutes $r_2(L')$ cycles of length $n=p^{\nu-c}$, since each complex
place downstairs has $n$ lifts. The absolute determinant of
$\gamma^{-1}-a$ on the permutation space is $(a^n-1)^{r_2(L')}$.
Its absolute determinant on the one-dimensional quotient is $a-1$.
The determinant on the lattice modulo roots of unity therefore gives
$|E'|(a-1)=(a^n-1)^{r_2(L')}$, as required.

The powered twisted norm induces a map
$W\to Y_F/(qY_F+Y_{F,\mathrm{tors}})$. Its image has size at most
$B=q^{r_2(L')}T_{2k+3}T_{2k+2}$ by choosing one representative for
each image element and applying the first count. Write
$a^{p^{\nu-c}}-1=p^{\nu+f}m_0$ with $p\nmid m_0$.
The norm identity kills the kernel by $p^{b+f}m_0$.
Indeed, write $g=\gamma^{-1}$, $n=|\gamma|$, $N_a=\sum_{j<n}a^jg^{n-j}$.
Then $(g-a)N_a=N_a(g-a)=(1-a^n)g$. If $p^bN_a y=qz+t$
with $t$ torsion, applying $g-a$ and dividing by $q$ modulo torsion
shows that $p^{b+f}m_0g y$ lies in $(g-a)Y_F+Y_{F,\mathrm{tors}}$.
This subgroup is invariant under $g$ and $g^{-1}$, proving the assertion.
On its $p$-primary subgroup this gives an exponent bound $p^{b+f}$.
The Selmer injection and the base-field Selmer count bound $|W[p]|$
by $p^g$. A finite abelian $p$-group killed by $p^{b+f}$, with
at most $p^g$ elements killed by $p$, has order at most $p^{(b+f)g}$.
Indeed multiplication by $p$ takes the subgroup killed by $p^{j+1}$
into the subgroup killed by $p^j$ with kernel of order at most $p^g$.
Induction from $j=0$ gives the bound. Applying it to the kernel gives
$p^{v_p(|W|)}\le Bp^{(b+f)g}$. In particular
\[
 h_\nu q^{r_2(L')}\le
 p^{v_p(|C'|)+v_p(|E'|)+c+f}
 =p^{v_p(|W|)+c+f}\le Bp^X.
\]
For the Selmer bound just used, lift $w\in W[p]$ to $u\in U$ and write
$pu=\delta y+t$ with $t$ torsion. If $\partial u=\delta_I d$, then
$\partial y=pd$ by injectivity of $\delta_I$. Also $gy-y$ vanishes
modulo $pY_F+$ torsion, since $p\mid a-1$. The class of the chosen
$y$ in that quotient distinguishes the $w$: equality of two such
classes gives $y-y'=pz+t'$, whence
$p(u-u'-\delta z)$ is torsion, as is $u-u'-\delta z$ itself.
To descend these classes, fix a generator $\rho$ of the roots of unity
of $F$. In multiplicative notation a representative $\alpha$ has
$\gamma\alpha/\alpha=v^p\rho^r$ for $0\le r<p$ after absorbing a
$p$th power. The pair $(r,N_{F/L'}v)$ takes at most $p^2$ values
because $(N_{F/L'}v)^p=(N_{F/L'}\rho)^{-r}$. Within a fiber the ratio
of two such $v$ has norm one. Hilbert 90 supplies a coboundary, so
after removal of a $p$th power the ratio of the two $\alpha$ descends
to $L'$. Outside $p$ the cyclotomic extension is unramified, so the
valuations of the descended element are divisible by $p$ there.
The resulting base-field classes modulo powers have size at most
$p^{2d'+\lfloor\log_2 h_{L'}\rfloor}$. Indeed valuations modulo $p$
at places over $p$ cost at most $p^{d'}$. With these fixed, taking
the $p$th root of the principal divisor costs at most
$|\mathrm{Cl}(L')[p]|\le p^{\lfloor\log_2 h_{L'}\rfloor}$.
The remaining units modulo $p$th powers cost at most $p^{d'}$
by the unit theorem. This gives the stated $p^g$.

Substituting the value of $B$ gives
\[
 h_\nu q^{r_2(L')}
 \ \le\ q^{r_2(L')}T_{2k+3}T_{2k+2}\,p^X.
\]
Cancel the positive common factor. All quantities on the right of
$h_\nu\le T_{2k+3}T_{2k+2}p^X$ are independent of $\nu$.
This uniformity allows passage back to the localization kernel:
restriction to the infinite cyclotomic tower is injective on that kernel,
and each finite set in its image is detected at a common finite level.
At that level the corresponding torsion class characters inject into
the functions counted above. Thus every such finite set has size at most
$T_{2k+3}T_{2k+2}p^X$. It follows that the kernel is finite and
\[
 s_L\log p\le s_{L'}\log p
 \le\log T_{2k+3}+\log T_{2k+2}+X\log p.
\]

For clarity on the constants, write $u=2ed'$ and $Y=\log(2D')$.
For the torsion estimate set $N=4e+6$ and $d_* = d'N$.
The integral face complex for $\mathrm{GL}_N(\mathcal O_{L'})$ is
exact in the range used to compute $H_j$. In dimensions up to $j$
its number of orbits is bounded by
$A_*=(2D'd_*)^{12(j+1)d_*^3N}$ and the orders of the stabilizers by
$s_*=3^{d_*N}$. The boundary from dimension $j+1$ is the alternating
sum of $j+2$ faces, hence has column $\ell^1$-norm at most $j+2$.
The equivariant chain estimate with these bounds gives
\[
 T_j\le (j+2)^{A_*}\,s_*^{\,jA_*s_*^j}.
\]
This estimate uses exactness of the augmented complex, orbit counting
for integral lattices, and the bounds for homology of finite stabilizers.
Taking logarithms gives
$A_*\log(j+2)+jA_*s_*^j\log s_*$. The numerical estimate for this
expression at $j\le2e+1$ gives
$\log T_j\le \exp(u^{C_0}Y)$ with $C_0=24$ in the two degrees above. The numerical
comparison for the layer gives $X\log p\le u^9Y$. Absorbing the three
terms and descending using the degree and discriminant bounds gives
$s_L\log p\le B_C(e,d,D)$ with $C=5(\max(C_0,9)+2)+11$.

\paragraph{From the localization kernel to the required order.}
Poitou--Tate duality and the local finiteness and duality inputs give
\[
 \operatorname{length}_{\mathbb Z_p} A
 \le s_L+\sum_{v\in S}\operatorname{length}_{\mathbb Z_p}H^0(G_{L_v},M).
\]
For each $v$ the cyclotomic image $U_v\subset\mathbb Z_p^\times$ has
index $I_v\le d$: it contains the inertia image, whose index is at most
$d$. The element $(1+p^2)^{I_v}$ lies in $U_v$.
An invariant of order $p^t$ forces
$p^t\mid (1+p^2)^{I_v e}-1$. The valuation of the latter is
$2+v_p(I_v)+v_p(e)$, also for $p=2$ since the base is $1\bmod 4$.
Consequently the invariants form a subgroup of a cyclic group of order
at most $p^{2+\lfloor\log_p d\rfloor+e}$. There are at most $d$ places
above $p$, so
\[
 \operatorname{length}_{\mathbb Z_p} A
 \le s_L+d(2+\lfloor\log_p d\rfloor+e).
\]
A finite-length $\mathbb Z_p$-module of length $t$ has order $p^t$,
by a composition series with successive quotients $\mathbb F_p$.
This proves finiteness. The root-of-unity hypothesis gives $p\le d+1$,
and $\lfloor\log_p d\rfloor\log p\le\log(d+1)$. Hence the extra term
in $\log|A|$ is at most $d(3+e)\log(d+1)$. For $e,d,D\ge1$ this is
at most $4d\log(2D)(2e)(2e+3)\log(d+1)$, which the final numerical
comparison absorbs into $B_{\max(C,0)+9}(e,d,D)$ together with $B_C$.
Following the exhibited witnesses gives $5(24+2)+11=141$ at the
localization-kernel step and $141+9=150$ here. The interface used for
the window only asks for the existence of one absolute exponent.

\paragraph{Scope of this reconstruction.}
This deduction uses the arithmetic and analytic inputs to the odd-degree
rank bound, the integral torsion estimate, the Selmer and class-group
comparisons, and the tower and duality theorems in their indicated
normalizations. Paragraph~\ref{kone:lines} gives the line transfer used in the cycle
argument. A standalone derivation still requires the integral lattice
and splitting results and the arithmetic and analytic inputs above. In the window argument
the same uniform bound is supplied to the case estimates together with the
separate rank-two Selmer input.

\paragraph{Use of the bound on the selected family.}\label{kone:consumer}
The window argument takes the existential pair $(C_K,\text{cohomology
bound at }C_K)$ once. It uses that same pair in the count of type I,
in the joint count of types II and III, and in the count of type V.
Type I also uses the rank-two Selmer estimate and the unramified
automorphy input. Types II and III use the symplectic input, while
type V uses the reduced unitary lifting input. The four counts and
the common pair yield the window estimate of
Hypothesis~\ref{w4:thm:vwin} in the source normalization.
At each selected index the window supplies finiteness of the odd
localization kernels and bounds the sum of their lengths, weighted by
$\log p$, by $\sqrt b$. Indeed the record identifies its finite
lengths with the cohomological ones at every $1\le e<t$; nonprime
indices contribute zero, and summation gives the required inequality.
With the local lattice comparison and equality of the small-level
ranks at that index this gives the twisted-place divisibility used
in the division estimate of Subsection~\ref{sec:frontier}.
These rank and automorphy hypotheses enter independently of the
cohomological order bound.
